% ******************************* Thesis Appendix A ****************************

\chapter{Software Requirements Specification (SRS)}
\label{appendix:srs}

\section{Introduction}

\subsection{Purpose \& Scope}
The University of San Carlos (USC) Clinic mandates a transition from a fragmented, paper-based operations model to a fully digital Health Information System (HIS). The USC-PIS is engineered to centralize and secure student health information by capturing comprehensive medical histories, treatment plans, medications, and dental records. The system scope encompasses the complete digital replacement of the Daily Treatment Records (DTR), consultation management, medical certificate issuance, health awareness campaigns, automated patient notifications, and rigorous feedback analytics.

\subsection{Design \& Implementation Constraints}
The "Modernizing the USC Clinic's Patient Information System" project is a web-based system developed under the timeframe and academic constraints of the CpE Design course sequence. As a student-led initiative, the following constraints and boundaries apply to the current deployment:
\begin{itemize}
    \item \textbf{Hardware \& Hosting Constraints:} The system relies on Heroku's basic tier hosting and PostgreSQL, deliberately excluding costly enterprise integrations like Single Sign-On (SSO).
    \item \textbf{Inventory Management:} Medication stock tracking is not included.
    \item \textbf{Billing Integration:} Advanced insurance processing and financial management are excluded.
    \item \textbf{Appointment Scheduling:} Consultations are strictly handled via a manual walk-in workflow as per clinic request.
    \item \textbf{Patient Medication Tracking:} This functionality is deferred to future iterations.
\end{itemize}
To maintain a strict boundary on the current system deployment, the following features are deliberately excluded from the current thesis scope and deferred to future enterprise phases.

\subsection{Definitions}
\begin{itemize}
    \item \textbf{PII/PHI (Personally Identifiable Information / Protected Health Information):} Sensitive patient data protected by the database architecture.
    \item \textbf{pgcrypto:} A native PostgreSQL extension integrated into the system for applying industry-standard, column-level symmetric encryption to sensitive data.
    \item \textbf{RBAC (Role-Based Access Control):} A security paradigm implemented via Django REST Framework (DRF) permission classes to restrict access dynamically.
    \item \textbf{FDI Notation:} The standard numerical tooth identification system (11-48) enforced via regex constraints within the Dental Records module.
\end{itemize}

\section{Overall Description}

\subsection{System Perspective}
The USC-PIS is built upon a modern, three-tier web architecture. The backend is powered by Django 5.0.2 and Django REST Framework, while the presentation layer operates on React 18 using a Vite build system and Material-UI components. The data layer relies on a dual-database environment featuring PostgreSQL for production and SQLite for local development, enforcing vendor-aware logic to prevent compatibility errors during testing. External enterprise integrations include Heroku for hosting, Cloudinary for scalable media storage, and the Gmail API utilizing OAuth 2.0 to bypass "less secure app" restrictions for clinical communications.

\subsection{User Classes \& Characteristics}
The system strictly enforces user authorizations across hierarchical classes:

\begin{table}[htbp]
\centering
\caption{User Classes \& Capabilities}
\begin{tabular}{|p{0.2\textwidth}|p{0.75\textwidth}|}
\hline
\textbf{Role} & \textbf{Access Level \& Capabilities} \\
\hline
\textbf{Admin} & Possesses full system access, including global RBAC adjustments, system health monitoring, email system administration, and report generation. \\
\hline
\textbf{Doctor / Dentist} & Endowed with full administrative and clinical access, only the Doctor role is uniquely authorized to approve medical certificates. Both can prescribe treatments and access all reporting/feedback analytics. \\
\hline
\textbf{Nurse} & Supervisory and data-entry capacity; authorized to cater to patients, record vital signs, input treatment plans, and draft medical certificates for doctor approval. \\
\hline
\textbf{Staff} & Manages clinic administration, patient intake, certificate issuance, report generation, and health campaign creation. All other clinical roles (Admin, Doctor, Dentist, and Nurse) possess the authorization to perform all tasks and access all functionalities available to this Staff role. \\
\hline
\textbf{Student / Faculty (Patient)} & Restricted to their own data; permitted only to read personal medical/dental timelines, submit visit-linked feedback, and view public health campaigns. \\
\hline
\end{tabular}
\end{table}

\subsection{Hardware Interfaces}
\begin{itemize}
    \item \textbf{Client Devices:} A standard desktop computer, laptop, or mobile device capable of rendering modern web applications.
    \item \textbf{Network:} A stable broadband internet connection is required for real-time data exchange with the cloud servers.
\end{itemize}

\subsection{Software Interfaces}
\begin{itemize}
    \item \textbf{User Interface (UI):} A web-based Graphical User Interface (GUI) built on React 18 and Material-UI, ensuring responsive, low-latency interactions across all devices.
    \item \textbf{Backend Framework:} Django 5.0.2 acts as the core application layer, providing secure RESTful API endpoints.
    \item \textbf{Database Management:} PostgreSQL handles relational data storage, utilizing the \texttt{pgcrypto} extension for data-at-rest encryption.
    \item \textbf{External APIs:} The system interfaces with the Gmail API (via OAuth 2.0) for automated communication, bypassing strict security blocks, and Cloudinary for scalable Content Delivery Network (CDN) file storage.
\end{itemize}

\section{Functional Requirements}

\begin{itemize}
    \item \textbf{Identity Management (authentication):} Secures access via the \texttt{/api/auth/} mount point. USC domain enforcement mandates that student registrations end in \texttt{@usc.edu.ph}. A \textbf{SafeList} registration feature maps authorized clinical emails to roles, bypassing standard MFA logic to automatically assign professional access.
    \item \textbf{Clinical Records \& Retooling (patients):} Data entry uses a unified record component (\texttt{MedicalRecord.jsx} and \texttt{Dental.jsx}). \textbf{Retooling Focus:} An advanced real-time filtering architecture processes queries via dynamic parameter filtering (\texttt{year\_level} and \texttt{academic\_year}). This logic calculates Semesters autonomously to filter the patients list accurately.
    \item \textbf{Health Awareness (health\_info):} Manages the lifecycle of clinical announcements across 13 defined categories (e.g., VACCINATION). Visibility is purely governed by \texttt{start\_date} and \texttt{end\_date} parameters; manual status toggling is deprecated.
    \item \textbf{Administrative Workflows (medical\_certificates):} Digitalizes the USC Form ACA-HSD-04F workflow. Nurses draft the certificate, but Doctors hold exclusive authority (\texttt{IsDoctorUser}) to approve the final document.
    \item \textbf{Feedback \& Notifications (feedback, notifications):} Feedback is linked uniquely to clinical visits via backend constraints, restricting users to one submission per visit. Creating a health record executes a signal that schedules a mandatory 24-hour follow-up alert.
    \item \textbf{Reports (reports):} Supports background processing logic (via Celery/Redis) to dynamically synthesize clinic reports, exporting data natively to PDF, Excel, and CSV without blocking HTTP threads.
\end{itemize}

\section{Data \& Security Requirements}

\begin{itemize}
    \item \textbf{Encryption-at-Rest:} Implements industry-level data encryption mapped to PostgreSQL's \texttt{pgcrypto} extension. Fields such as \texttt{illness\_enc} and \texttt{diagnosis\_enc} bypass plain-text memory and are converted strictly to \texttt{BinaryField} bytea hashes using backend signals upon database commits.
    \item \textbf{Access Control:} Guarded globally by RBAC middleware using DRF \texttt{PermissionClasses}. Unauthenticated traffic is blocked by default. Files accessed under \texttt{/api/files/patient-documents/} are proxied strictly through authenticated backend streaming, blocking public CDN bypass attempts.
\end{itemize}

\subsection{System Accountability and Audit Logging}
To ensure data integrity, operational transparency, and adherence to healthcare compliance standards, the USC-PIS implements a comprehensive tracking and logging architecture across its core modules.

\begin{itemize}
    \item \textbf{Database Audit Trails:} The system enforces strict data accountability by embedding structural audit trails within the database schema. All critical entities, including patient profiles, medical records, and dental records, are automatically equipped with \texttt{created\_at} and \texttt{updated\_at} timestamp fields. This guarantees a precise, immutable history of when clinical data was first entered and subsequently modified.
    \item \textbf{System Activity and Notification Logging:} To monitor administrative oversight and system interactions, the USC-PIS maintains a persistent log of all dispatched email and in-app communications. Administrators have access to detailed "System Activity Logs" via the Email Administration dashboard, which provides timestamped records of backend email events and error diagnostics. Crucially, this creates an explicit audit trail that tracks delivery statuses and identifies exactly which administrator triggered specific manual actions or system overrides.
    \item \textbf{API Log Access and Compliance:} All user actions and email administration tasks are systematically logged to prevent unauthorized system manipulation and to ensure security. Furthermore, any manual email dispatches executed by the staff are explicitly documented for audit purposes. Authorized administrators can programmatically review these communication audit trails via the \texttt{/api/notifications/logs/} endpoint.
\end{itemize}

\section{Non-Functional Requirements \& Future Roadmap}

\subsection{Performance \& Security Benchmarks}
System performance benchmarking explicitly addressed and resolved prior inefficiencies:

\begin{table}[htbp]
\centering
\caption{System Benchmarks}
\begin{tabular}{|p{0.2\textwidth}|p{0.6\textwidth}|p{0.15\textwidth}|}
\hline
\textbf{Requirement Area} & \textbf{Verified Benchmark} & \textbf{Audit Test ID} \\
\hline
\textbf{Search Latency} & Patient searches execute in \textbf{127.81ms} (Target: <500ms). & PT-04 \\
\hline
\textbf{PDF Generation} & Medical Certificate PDFs render in \textbf{122.93ms}. & PT-01 \\
\hline
\textbf{Concurrency} & Processes 20 concurrent requests in \textbf{0.66 seconds} (0\% failure). & PT-02 \\
\hline
\textbf{File Security} & \textbf{100\% rejection rate} against executable scripts (.EXE, .JS). & UT-06 \\
\hline
\end{tabular}
\end{table}

\subsection{Data Resilience}
\textbf{Backup and Recovery:} The system is equipped with an automated data resilience module that allows authorized administrators to export the entire structured database, including patient identifiers, medical histories, and user accounts, into a single, highly compressed JSON file. This ensures business continuity and provides a rapid recovery point in the event of environment-level failures.

\subsection{Future Roadmap}
While current features solve immediate inefficiencies, the system's scalable architecture supports future enhancements:
\begin{itemize}
    \item \textbf{Mobile Progressive Web App (PWA):} Deploying a mobile application wrapper with offline capabilities and push notifications.
    \item \textbf{Telemedicine Platform:} Expanding healthcare delivery to include video consultations and remote monitoring.
    \item \textbf{AI/ML Integration:} Implementing machine learning models for predictive health analytics and diagnostic assistance.
\end{itemize}
